\documentclass{article}
\usepackage{amsmath}
\usepackage{amssymb}
\usepackage{physics}
\usepackage{bm}

\title{Quantum Mechanics 3 -- Question 2, Task 7 Ext.}
\author{Kieran Ody}
\date{August 2026}

\begin{document}

\maketitle

\section*{2. The Particle in a One-Dimensional Infinite Potential Well -- Including Task 7 Extension}

The infinite potential well (or ``particle in a box'') is one of the simplest exactly solvable quantum systems. Although idealised, it demonstrates several fundamental concepts of quantum mechanics, including wavefunctions, quantised energy levels and Heisenberg's Uncertainty Principle.

\begin{equation*}
\psi(x,t)
\end{equation*}

describes the wavefunction of a particle of mass $m$ in a `box' of potential energy V.

\vspace{1\baselineskip}

The potential is defined as

\begin{equation*}
V(x)=
\begin{cases}
0, & 0<x<a,\\
\infty, & \text{otherwise}.
\end{cases}
\end{equation*}


\subsection*{2(i) Standing Waves}

The particle is confined to the region $0\leq x\leq a$ by an
infinite potential outside the box. Therefore, the wavefunction
must satisfy the boundary conditions

\begin{equation*}
\psi(0,t)=\psi(a,t)=0
\end{equation*}

At the left-hand boundary, $x=0$, the sine term becomes

\begin{equation*}
\sin(0)=0,
\end{equation*}

so that $\psi(0,t)=0$. At the right-hand boundary, $x=a$,

\begin{equation*}
\sin\left(\frac{n\pi a}{a}\right) = \sin(n\pi) = 0
\end{equation*}

This means the points $x=0$ and $x=a$ are nodes of a standing wave.

\vspace{1\baselineskip}

The condition above is satisfied when $n$ is an integer. The value
$n=0$ would result in $\psi(x,t)=0$ everywhere and therefore cannot
represent a particle, so the allowed values are

\begin{equation*}
\boxed{n=1,2,3,\ldots}
\end{equation*}

Thus, $n$ is a positive integer and represents the quantum number
associated with the standing-wave state.

\begin{equation*}
\cos(\omega t)
\end{equation*}

where $\omega = 2\pi f$, $f=c/\lambda$, describes the wavefunction's sinusoidal time dependence. At a fixed position $x$, the wavefunction therefore oscillates with angular frequency $\omega$:

\begin{equation*}
\psi(x,t)\propto\cos(\omega t)
\end{equation*}

The angular frequency is related to the energy of the state by

\begin{equation*}
E=\hbar\omega
\end{equation*}

\vspace{1\baselineskip}

A suitable form for the wavefunction is therefore

\begin{equation*}
\boxed{\psi(x,t) = A_n\sin\left(\frac{n\pi x}{a}\right)\cos(\omega t)}
\end{equation*}

$A_n$ is the normalisation factor. The sine term describes the
spatial standing-wave pattern, while the cosine term
describes its time dependence.


\subsection*{2(ii) A Solution to the Schrödinger Equation}

The time-independent Schrödinger equation is

\begin{equation*}
-\frac{\hbar^2}{2m} \frac{\partial^2\psi}{\partial x^2} + V\psi = E\psi
\end{equation*}

Inside the box, the potential is zero, and the time-independent Schrödinger equation is therefore

\begin{equation*}
-\frac{\hbar^2}{2m} \frac{\partial^2\psi}{\partial x^2} = E\psi
\end{equation*}

Consider the proposed wavefunction

\begin{equation*}
\psi(x,t) = A_n\sin\left(\frac{n\pi x}{a}\right)\cos(\omega t)
\end{equation*}

Differentiating twice with respect to $x$ gives

\begin{equation*}
\frac{\partial^2\psi}{\partial x^2} = -A_n \left(\frac{n\pi}{a}\right)^2 \sin\left(\frac{n\pi x}{a}\right) \cos(\omega t)
= -\left(\frac{n\pi}{a}\right)^2\psi
\end{equation*}

Substituting this into the Schrödinger equation gives

\begin{equation*}
-\frac{\hbar^2}{2m} \left[-\left(\frac{n\pi}{a}\right)^2\psi\right] = E\psi
\end{equation*}

Hence,

\begin{equation*}
\boxed{E(n,a) = \frac{n^2\pi^2\hbar^2}{2ma^2}}
\end{equation*}

For a proton confined to a box of width

\begin{equation*}
a = 0.47\times10^{-15}\,\mathrm{m},
\end{equation*}

the ground-state energy ($n=1$) is approximately

\begin{equation*}
E_1\approx927\,\mathrm{MeV}
\end{equation*}

Since the energy varies with $n^2$, the general energy is

\begin{equation*}
\boxed{E_n\approx927n^2\,\mathrm{MeV}}
\end{equation*}

The proton rest mass energy is

\begin{equation*}
m_pc^2\approx938\,\mathrm{MeV}
\end{equation*}

Consequently,

\begin{equation*}
\frac{E_1}{m_pc^2} \approx \frac{927}{938} \approx 0.988
\end{equation*}

The ground-state energy is approximately $98.8\%$ of the proton's
rest mass energy, $m_pc^2$. The non-relativistic Schrödinger equation
assumes that the kinetic energy is much smaller than the rest energy,
so that

\begin{equation*}
K_e \approx \frac{p^2}{2m}.
\end{equation*}

However, in this case, the predicted confinement energy is comparable
to $m_pc^2$. This indicates that the proton would have a sufficiently
large momentum for relativistic effects to become important. The non-relativistic particle-in-a-box model no longer provides an accurate description of the proton in such a small confinement region. A relativistic approach would instead be required.


\subsection*{2(iii) Time Dependence and Normalisation}

The time-dependent factor $\cos(\omega t)$ can be replaced by the
complex exponential $e^{-iEt/\hbar}$. This follows from the
relationship between energy and angular frequency,

\begin{equation*}
E=\hbar\omega
\end{equation*}

Thus, the wavefunction can be written as

\begin{equation*}
\psi_n(x,t) = A_n \sin\left(\frac{n\pi x}{a}\right) e^{-iE_nt/\hbar}
\end{equation*}

The Born interpretation states that $|\psi(x,t)|^2$ represents the
probability density of finding the particle at position $x$. Since
the particle must be found somewhere within the box, the total
probability must be total certainty:

\begin{equation*}
\int_0^a |\psi_n(x,t)|^2\,dx=1
\end{equation*}

Taking the modulus squared of the wavefunction gives

\begin{equation*}
|\psi_n(x,t)|^2 = A_n^2 \sin^2\left(\frac{n\pi x}{a}\right) \left|e^{ iE_nt/\hbar}\right|^2
\end{equation*}

Since

\begin{equation*}
\left|e^{-iE_nt/\hbar}\right|^2 = 1,
\end{equation*}

the probability density is independent of time:

\begin{equation*}
|\psi_n(x,t)|^2 = A_n^2 \sin^2\left(\frac{n\pi x}{a}\right)
\end{equation*}

Normalisation therefore requires

\begin{equation*}
A_n^2 \int_0^a \sin^2\left(\frac{n\pi x}{a}\right)\,dx = 1
\end{equation*}

Using

\begin{equation*}
\int_0^a \sin^2\left(\frac{n\pi x}{a}\right)\,dx = \frac{a}{2},
\end{equation*}

we obtain

\begin{equation*}
A_n^2\frac{a}{2}=1
\end{equation*}

Therefore,

\begin{equation*}
\boxed{A_n=\sqrt{\frac{2}{a}}}
\end{equation*}

The normalised wavefunction is

\begin{equation*}
\boxed{\psi_n(x,t) = \sqrt{\frac{2}{a}} \sin\left(\frac{n\pi x}{a}\right) e^{-iE_nt/\hbar}}
\end{equation*}


\subsection*{2(iv) Expected Values of Momentum}

The momentum operator in one dimension is

\begin{equation*}
\hat{p}=-i\hbar\frac{\partial}{\partial x}
\end{equation*}

The expected value of momentum is given by

\begin{equation*}
\langle p\rangle = \int_0^a \psi^*\hat{p}\psi\,dx
\end{equation*}

For the normalised wavefunction

\begin{equation*}
\psi_n(x,t) = \sqrt{\frac{2}{a}} \sin\left(\frac{n\pi x}{a}\right) e^{-iE_nt/\hbar},
\end{equation*}

differentiation with respect to $x$ gives

\begin{equation*}
\frac{\partial\psi_n}{\partial x} = \sqrt{\frac{2}{a}} \frac{n\pi}{a} \cos\left(\frac{n\pi x}{a}\right) e^{-iE_nt/\hbar}
\end{equation*}

Therefore,

\begin{equation*}
\langle p\rangle = -i\hbar\frac{2n\pi}{a^2} \int_0^a \sin\left(\frac{n\pi x} {a}\right) \cos\left(\frac{n\pi x}{a}\right) \,dx
\end{equation*}

The integral evaluates to zero, giving

\begin{equation*}
\boxed{\langle p\rangle=0}
\end{equation*}

Physically, this occurs because the standing wave contains equal
components travelling in opposite directions, corresponding to equal
and opposite momenta.

\vspace{1\baselineskip}

The square of the momentum operator is

\begin{equation*}
\hat{p}^2 = -\hbar^2\frac{\partial^2}{\partial x^2}
\end{equation*}

From the second derivative of the wavefunction,

\begin{equation*}
\frac{\partial^2\psi_n}{\partial x^2} = -\left(\frac{n\pi}{a}\right)^2\psi_n
\end{equation*}

Substituting,

\begin{equation*}
\hat{p}^2\psi_n = \frac{n^2\pi^2\hbar^2}{a^2}\psi_n
\end{equation*}

Taking the expected value gives

\begin{equation*}
\langle p^2\rangle = \frac{n^2\pi^2\hbar^2}{a^2} \int_0^a|\psi_n|^2\,dx
\end{equation*}

Since the wavefunction is normalised,

\begin{equation*}
\int_0^a|\psi_n|^2\,dx=1,
\end{equation*}

and hence

\begin{equation*}
\boxed{ \langle p^2\rangle = \frac{n^2\pi^2\hbar^2}{a^2}}
\end{equation*}


\subsection*{2(v)(a) Expected Values of $x$}

The expected value of x is given by

\begin{equation*}
\langle x\rangle = \int_0^a x|\psi_n|^2\,dx = \frac{2}{a} \int_0^a x\sin^2\left(\frac{n\pi x}{a}\right)\,dx
\end{equation*}

Quoting the integral identity:

\begin{equation*}
\int x \sin^2(\alpha x) \, dx = \frac{1}{4}x^2 - \frac{x\sin(2\alpha x)}{4\alpha} - \frac{\cos(2\alpha x)}{8\alpha^2}
\end{equation*}

Let

\begin{equation*}
\alpha = \frac{n\pi}{a}
\end{equation*}

Evaluating the definite integral from $x = 0$ to $x = a$ (noting that $2\alpha a = 2n\pi$, so $\sin(2n\pi) = 0$ and $\cos(2n\pi) = 1$):

\begin{align}
\int_0^a x \sin^2(\alpha x) \, dx &= \left[ \frac{1}{4}x^2 - \frac{x\sin(2\alpha x)}{4\alpha} - \frac{\cos(2\alpha x)}{8\alpha^2} \right]_0^a \\
&= \left( \frac{1}{4}a^2 - 0 - \frac{1}{8\alpha^2} \right) - \left( 0 - 0 - \frac{1}{8\alpha^2} \right) \\
&= \frac{1}{4}a^2
\end{align}

Thus:
\begin{equation*}
\boxed{\langle x \rangle = \frac{2}{a} \left( \frac{1}{4}a^2 \right) = \frac{1}{2}a}
\end{equation*}


\subsection*{2(v)(b) Calculation of $\langle x^2 \rangle$}

The expected value $\langle x^2 \rangle$ is:

\begin{equation*}
\langle x^2 \rangle = \int_0^a x^2 |\psi_n(x)|^2 \, dx = \frac{2}{a} \int_0^a x^2 \sin^2(\alpha x) \, dx
\end{equation*}

Quoting the integral identity:

\begin{equation*}
\int x^2 \sin^2(\alpha x) \, dx = \frac{1}{6}x^3 - \frac{x \cos(2\alpha x)}{4\alpha^2} - \frac{(2\alpha^2 x^2 - 1)\sin(2\alpha x)}{8\alpha^3}
\end{equation*}

Let

\begin{equation*}
\alpha = \frac{n\pi}{a}
\end{equation*}

Evaluating from $x = 0$ to $x = a$:

\begin{align*}
\int_0^a x^2 \sin^2(\alpha x) \, dx &= \left[ \frac{1}{6}x^3 - \frac{x \cos(2\alpha x)}{4\alpha^2} - \frac{(2\alpha^2 x^2 - 1)\sin(2\alpha x)}{8\alpha^3} \right]_0^a \\
&= \left( \frac{1}{6}a^3 - \frac{a}{4\alpha^2} - 0 \right) - (0 - 0 - 0) \\
&= \frac{a^3}{6} - \frac{a}{4\alpha^2}
\end{align*}

Substituting $\alpha = \frac{n\pi}{a}$:

\begin{equation*}
\int_0^a x^2 \sin^2(\alpha x) \, dx = \frac{a^3}{6} - \frac{a^3}{4n^2\pi^2}
\end{equation*}

And

\begin{align*}
\langle x^2 \rangle &= \frac{2}{a} \left( \frac{a^3}{6} - \frac{a^3}{4n^2\pi^2} \right) \\
&= \frac{a^2}{3} - \frac{a^2}{2n^2\pi^2} \\
&= \frac{1}{3}a^2 \left( 1 - \frac{3}{2n^2\pi^2} \right)
\end{align*}


\subsection*{2(vi)(a) Calculation of $\Delta x$}

The uncertainty in position $\Delta x$ is defined by:

\begin{equation*}
(\Delta x)^2 = \langle x^2 \rangle - \langle x \rangle^2
\end{equation*}

Using the previously derived results of $\langle x \rangle$ and $\langle x^2 \rangle$:

\begin{align*}
(\Delta x)^2 &= \left( \frac{a^2}{3} - \frac{a^2}{2n^2\pi^2} \right) - \left( \frac{a}{2} \right)^2 \\
&= \left( \frac{a^2}{3} - \frac{a^2}{4} \right) - \frac{a^2}{2n^2\pi^2} \\
&= \frac{a^2}{12} - \frac{a^2}{2n^2\pi^2}
\end{align*}

Factoring out $\frac{a^2}{4n^2\pi^2}$:
\begin{align*}
(\Delta x)^2 &= \frac{a^2}{4n^2\pi^2} \left( \frac{4n^2\pi^2}{12} - 2 \right) \\
&= \frac{a^2}{4n^2\pi^2} \left( \frac{n^2\pi^2}{3} - 2 \right)
\end{align*}

Taking the square root yields
\begin{equation*}
\Delta x = \frac{a}{2n\pi} \sqrt{\frac{n^2\pi^2}{3} - 2}
\end{equation*}

\subsection*{2(vi)(b) Calculation of $\Delta p$}

The uncertainty in momentum $\Delta p$ is defined by:
\begin{equation*}
(\Delta p)^2 = \langle p^2 \rangle - \langle p \rangle^2
\end{equation*}

Using the previously derived results of $\langle p \rangle$ and $\langle p^2 \rangle$:
\begin{equation*}
\Delta p = \sqrt{\frac{n^2\pi^2\hbar^2}{a^2} - 0} = \frac{n\pi\hbar}{a}
\end{equation*}

\subsection*{2(vi)(c) Product of Uncertainties $\Delta x \Delta p$}

Multiplying $\Delta x$ and $\Delta p$:
\begin{align*}
\Delta x \Delta p &= \left( \frac{a}{2n\pi} \sqrt{\frac{n^2\pi^2}{3} - 2} \right) \left( \frac{n\pi\hbar}{a} \right) \\
&= \frac{1}{2}\hbar \sqrt{\frac{n^2\pi^2}{3} - 2}
\end{align*}

\subsection*{2(vi)(d) Verification of the Uncertainty Principle}

Heisenberg's Uncertainty Principle requires that:
\begin{equation*}
\Delta x \Delta p \ge \frac{\hbar}{2}
\end{equation*}

For our result to satisfy this inequality, we need:
\begin{equation*}
\sqrt{\frac{n^2\pi^2}{3} - 2} \ge 1
\end{equation*}

Testing for the ground state ($n = 1$):
\begin{equation*}
\frac{(1)^2\pi^2}{3} - 2 \approx 1.2899 > 1
\end{equation*}

Taking the square root:
\begin{equation*}
\sqrt{1.2899} \approx 1.1357 > 1
\end{equation*}

Since $\sqrt{\frac{n^2\pi^2}{3} - 2}$ increases as $n$ increases, the factor is strictly greater than 1 for all quantum numbers $n$, therefore,
\begin{equation*}
\Delta x \Delta p > \frac{\hbar}{2} \quad \text{for all } n = 1, 2, 3, \dots
\end{equation*}

This verifies that the particle in a box always satisfies Heisenberg's Uncertainty Principle.

\end{document}
